\section{E4: First-submission financial code generation}
\label{sec:codegen}
Can an open code model write a correct allocation or signal rule from a specification alone,
and does placing the library's documentation in context change that? We froze a paired study
before model inference. Qwen2.5-Coder 7B and 14B each wrote one program for four rules
(diagonal risk parity by inverse volatility, positive cross-sectional momentum, a 5/20
moving-average crossover, and volatility-targeted momentum), three data seeds, and two arms:
\CGPlanned{} programs. Both arms received identical raw quotes with split jumps, cumulative
split factors known on their event dates, the mathematical rule, the one-bar timing rule, one
response of at most 2,048 tokens, and no execution feedback. The library arm additionally
received the actual catalog output and interface and guard documentation (mean prompt
\CGFourteenLibraryInTokens{} versus \CGFourteenRawInTokens{} tokens). Programs ran in the
Landlock/seccomp worker. An independent NumPy grader, shared by both arms, requires
execution, agreement with the reference within $10^{-8}$, unchanged positions through
session $k$ when prices at and after $k$ are perturbed ($k=32,64$), and identical output for
the economically equivalent split-adjusted input. A program is valid only if it passes all.

\begin{table*}[t]
\centering\small
\setlength{\tabcolsep}{5pt}
\begin{tabular}{llcccccccc}
\toprule
& & \multicolumn{2}{c}{Executed} & \multicolumn{4}{c}{Post-hoc extraction} &
\multicolumn{2}{c}{Mean per response} \\
\cmidrule(lr){3-4}\cmidrule(lr){5-8}\cmidrule(lr){9-10}
Model & Arm & Official & Post-hoc & Correct & Split fail & Timing fail & Valid &
Output tok. & Gen.\ s \\
\midrule
7B & Raw & \CGSevenRawOfficialExec/12 & \CGSevenRawPostExec/12 & \CGSevenRawPostNumerical &
\CGSevenRawPostSplitFail & \CGSevenRawPostTimingFail & \CGSevenRawPostValid/12 &
\CGSevenRawOutTokens & \CGSevenRawGenSeconds \\
7B & Library & \CGSevenLibraryOfficialExec/12 & \CGSevenLibraryPostExec/12 &
\CGSevenLibraryPostNumerical & \CGSevenLibraryPostSplitFail & \CGSevenLibraryPostTimingFail &
\CGSevenLibraryPostValid/12 & \CGSevenLibraryOutTokens & \CGSevenLibraryGenSeconds \\
14B & Raw & \CGFourteenRawOfficialExec/12 & \CGFourteenRawPostExec/12 &
\CGFourteenRawPostNumerical & \CGFourteenRawPostSplitFail & \CGFourteenRawPostTimingFail &
\CGFourteenRawPostValid/12 & \CGFourteenRawOutTokens & \CGFourteenRawGenSeconds \\
14B & Library & \CGFourteenLibraryOfficialExec/12 & \CGFourteenLibraryPostExec/12 &
\CGFourteenLibraryPostNumerical & \CGFourteenLibraryPostSplitFail &
\CGFourteenLibraryPostTimingFail & \CGFourteenLibraryPostValid/12 &
\CGFourteenLibraryOutTokens & \CGFourteenLibraryGenSeconds \\
\bottomrule
\end{tabular}
\caption{Paired first-submission code generation. Under the pre-registered parser no program
is valid in any cell (\CGOfficialValid/\CGPlanned). Post-hoc columns take the first complete
code fence; failure counts are among executed programs. No library-arm program called a
library algorithm or guard.}
\label{tab:codegen}
\end{table*}

\paragraph{Pre-registered outcome.}
No program was valid in either arm (\CGOfficialValid/\CGPlanned). The protocol asked for
Python only, optionally inside one fence. In all, \CGFormatFailures{} responses appended an
explanation after the fence (\CGSevenTrailingProse{} from 7B and \CGFourteenTrailingProse{} from
14B), so the frozen parser compiled prose and recorded a \texttt{SyntaxError}. The
\CGOfficialExecuted{} executed programs all failed numerical agreement and split invariance.
No response was truncated. None of the \CGLibraryArm{} library-arm responses imported or
mentioned the library, and no algorithm or guard call was traced. The comparison therefore
measures documentation exposure without uptake; it does not estimate the effect of using
\textsc{fin-skills}.

\paragraph{Post-hoc extraction and failure modes.}
To separate formatting from computation, we froze an arm-blind extraction of the first
complete fence by hash before execution and reran the same worker, sandbox, and grader on
another Linux host; it reproduced all \CGEnvReproduced{} official executions exactly.
Then \CGPostExecuted{} programs execute, but only \CGPostValid{} of \CGPlanned{} is valid.
The failures are silent in the sense that matters for research: \CGRecumulated{} of 24 14B
programs, in both arms, applied a cumulative product to split factors that the prompt defines
as already cumulative, and \CGPostSplitFailures{} executed programs change their positions
when the same prices are presented split-adjusted. In \CGPostTimingFailures{} executed
programs a position moves when prices at or after its own session change; \CGSameSession{}
use the same-session close and \CGStrictlyFuture{} use later data. The 7B programs that still fail
raise shape, index, or pandas-API errors.

\paragraph{What the library guard observed.}
Run post hoc on each executed program, the library's causality guard flagged all
\CGGuardCaught{} programs with strictly-future dependence and none of the \CGTimingClean{}
programs that passed the grader's timing probes. It passed the \CGSameSessionGuardPassed{}
same-session violations: the guard asks whether values before $k$ depend on data at or after
$k$, which is feature causality, whereas a position must also exclude session $k$ itself.
The execution-lag contract therefore needs its own check rather than an inference from a
passed causality test. Together with E2, these results show that raw first submissions are
unreliable in ways that still execute, and that optional documentation is not used; they
motivate controller-enforced library calls and guards, which remain to be tested on these
tasks.
